% config/hyperref.tex — 超链接、PDF 书签与交叉引用
% ===================================================

% hyperref 必须在大多数包之后加载
\usepackage{hyperref}
\usepackage{bookmark}   % 更好的 PDF 书签，消除 rerun 警告
\renewcommand{\UrlFont}{\TTTFont}

\ifprintmode
  \hypersetup{
    colorlinks         = false,
    pdfborder          = {0 0 0},
    hypertexnames      = false,
    bookmarksnumbered  = true,
    bookmarksopen      = true,
    bookmarksopenlevel = 2,
    pdfstartview       = FitH,
    pdftitle           = {高级宏观经济学笔记},
    pdfauthor          = {作者},
    pdfsubject         = {Advanced Macroeconomics},
  }
\else
  \hypersetup{
    colorlinks         = true,
    linkcolor          = fafugreen,
    citecolor          = fafugreen,
    urlcolor           = fafugreen,
    hypertexnames      = false,
    bookmarksnumbered  = true,
    bookmarksopen      = true,
    bookmarksopenlevel = 2,
    pdfstartview       = FitH,
    pdftitle           = {高级宏观经济学笔记},
    pdfauthor          = {作者},
    pdfsubject         = {Advanced Macroeconomics},
  }
\fi

% cleveref 必须在 hyperref 之后加载
\usepackage[nameinlink]{cleveref}

% ----- 数学公式编号：章-序号（如 1-1） -----
\makeatletter
\renewcommand{\theequation}{\thechapter-\arabic{equation}}
\renewcommand{\theHequation}{chap.\thechapter.eq.\arabic{equation}}
\renewcommand{\theHfigure}{chap.\thechapter.fig.\arabic{figure}}
\renewcommand{\theHtable}{chap.\thechapter.tab.\arabic{table}}
\makeatother

% ----- cleveref 中文名称与格式 -----
\crefformat{chapter}{第~#2#1#3~章}
\Crefformat{chapter}{第~#2#1#3~章}
\crefrangeformat{chapter}{第~#3#1#4~章至第~#5#2#6~章}

\crefformat{section}{第~#2#1#3~节}
\Crefformat{section}{第~#2#1#3~节}
\crefrangeformat{section}{第~#3#1#4~节至第~#5#2#6~节}
\crefformat{subsection}{第~#2#1#3~小节}
\Crefformat{subsection}{第~#2#1#3~小节}
\crefrangeformat{subsection}{第~#3#1#4~小节至第~#5#2#6~小节}

\crefformat{equation}{式~(#2#1#3)}
\Crefformat{equation}{式~(#2#1#3)}
\crefrangeformat{equation}{式~(#3#1#4)\textendash (#5#2#6)}

\crefname{figure}{图}{图}
\Crefname{figure}{图}{图}

\crefname{table}{表}{表}
\Crefname{table}{表}{表}

% ----- \ref 显示“编号 + 标题名” -----
% 对 ch:/sec:/subsec: 前缀标签生效；其他标签保持标准 \ref 文本。
\let\AdvMacroRawRef\ref
\newcommand{\AdvMacroRefWithTitle}[1]{%
  \IfBeginWith{#1}{ch:}
    {第~\AdvMacroRawRef{#1}~章《\nameref*{#1}》}
    {%
      \IfBeginWith{#1}{part:}
        {第~\AdvMacroRawRef{#1}~部分《\nameref*{#1}》}
        {%
      \IfBeginWith{#1}{sec:}
        {第~\AdvMacroRawRef{#1}~节《\nameref*{#1}》}
        {%
          \IfBeginWith{#1}{subsec:}
            {第~\AdvMacroRawRef{#1}~小节《\nameref*{#1}》}
            {%
              \IfBeginWith{#1}{fig:}
                {图~\AdvMacroRawRef{#1}}
                {%
                  \IfBeginWith{#1}{tab:}
                    {表~\AdvMacroRawRef{#1}}
                    {%
                      \IfBeginWith{#1}{eq:}
                        {(\AdvMacroRawRef{#1})}
                        {\AdvMacroRawRef{#1}}%
                    }%
                }%
            }%
        }%
        }%
    }%
}
\AtBeginDocument{%
  \let\AdvMacroRawRef\ref
  \let\ref\AdvMacroRefWithTitle
}

% ----- tcolorbox 定理类计数器映射 -----
\makeatletter
\crefalias{tcb@cnt@definitionbox}{definition}
\crefalias{tcb@cnt@theorembox}{theorem}
\crefalias{tcb@cnt@lemmabox}{lemma}
\crefalias{tcb@cnt@propositionbox}{proposition}
\crefalias{tcb@cnt@corollarybox}{corollary}
\crefalias{tcb@cnt@assumptionbox}{assumption}
\crefalias{tcb@cnt@examplebox}{example}
\makeatother

\crefname{definition}{定义}{定义}
\Crefname{definition}{定义}{定义}

\crefname{theorem}{定理}{定理}
\Crefname{theorem}{定理}{定理}

\crefname{lemma}{引理}{引理}
\Crefname{lemma}{引理}{引理}

\crefname{proposition}{命题}{命题}
\Crefname{proposition}{命题}{命题}

\crefname{corollary}{推论}{推论}
\Crefname{corollary}{推论}{推论}

\crefname{assumption}{假设}{假设}
\Crefname{assumption}{假设}{假设}

\crefname{example}{例}{例}
\Crefname{example}{例}{例}
